\section{Limitations}
\label{sec:limitations}

No detector is checked against human labels. Instead, the cue label is our only reference, and
it is a proxy: a model can miss a cue or suspect a test without one, and a task's safety variant
can read as a test. Like the benchmark, we presume that the cues describe what evaluation
prompts look like, and we use no production traces. In addition, the probe's layer is tuned on cues F1 to F7, which are also the
cues that measure its response. Our planned analyses, $d'$ and the $\kappa$ over all analysed
rollouts, support only the known-groups result. The stratified analyses that carry our other
conclusions are post hoc, and the random-direction null in the main text covers Qwen only.
Moreover, most cued Qwen
rollouts carry several cues (\cv{design/expB} multi-cue of \cv{qwen/cuedrollouts} cued
rollouts). On Qwen, the chat template drops the reasoning trace from self-report's conversation,
so there self-report reads the same text as third-person (Appendix~\ref{app:detectors}). The
probe result holds only for one contrastive direction from \cv{design/pairs} prompt pairs. Other
probes could behave differently: a logistic-regression probe trained on the cue labels reaches an
AUROC of \cv{qwen/logistic/auroc} (Figure~\ref{fig:probe-null}). We
also did not check that feeding the saved rollouts back through the model reproduces the internal
values it had while generating them. Finally, the scope is two models, one
benchmark and \cv{design/reps} repetitions per prompt. Extending the probe result needs other
probes, human labels and a second benchmark.
